\documentclass[11pt,a4paper]{article}
% preamble.tex — estilo común de todos los apuntes Froth.
% Cada apunte hace % preamble.tex — estilo común de todos los apuntes Froth.
% Cada apunte hace % preamble.tex — estilo común de todos los apuntes Froth.
% Cada apunte hace \input{preamble} justo después de \documentclass.
\usepackage[margin=2.4cm]{geometry}
\usepackage{xcolor}
\definecolor{litblue}{RGB}{31,79,121}
\definecolor{litgreen}{RGB}{27,94,32}
\definecolor{litorange}{RGB}{230,81,0}
\usepackage{parskip}
\usepackage{titlesec}
\titleformat{\section}{\Large\bfseries\color{litblue}}{\thesection}{0.6em}{}
\titleformat{\subsection}{\large\bfseries\color{litblue}}{\thesubsection}{0.6em}{}
\usepackage{tcolorbox}
\usepackage{fancyhdr}
\pagestyle{fancy}
\fancyhf{}
\fancyhead[L]{\small\color{litblue}Froth — Apuntes de ML}
\fancyhead[R]{\small\thepage}
\renewcommand{\headrulewidth}{0.4pt}
\usepackage[colorlinks=true,linkcolor=litblue,urlcolor=litblue]{hyperref}

% Cabecera de cada apunte: \cabecera{numero}{titulo}
\newcommand{\cabecera}[2]{%
  \begin{tcolorbox}[colback=litblue,coltext=white,colframe=litblue,arc=2mm]
    {\Large\bfseries Apunte #1 — #2}\\[3pt]
    {\small Proyecto Froth · aprender Machine Learning construyendo}
  \end{tcolorbox}\vspace{0.8em}}

% Cajas temáticas
\newtcolorbox{concepto}[1]{colback=blue!4,colframe=litblue,title={Concepto: #1},arc=1.5mm}
\newtcolorbox{practica}{colback=green!5,colframe=litgreen,title={Pruébalo tú (teclado en mano)},arc=1.5mm}
\newtcolorbox{ojo}{colback=orange!8,colframe=litorange,title={Ojo / error típico},arc=1.5mm}
\newtcolorbox{resumen}{colback=gray!8,colframe=gray!60,title={En una frase},arc=1.5mm}

\newcommand{\codigo}[1]{\texttt{#1}}
 justo después de \documentclass.
\usepackage[margin=2.4cm]{geometry}
\usepackage{xcolor}
\definecolor{litblue}{RGB}{31,79,121}
\definecolor{litgreen}{RGB}{27,94,32}
\definecolor{litorange}{RGB}{230,81,0}
\usepackage{parskip}
\usepackage{titlesec}
\titleformat{\section}{\Large\bfseries\color{litblue}}{\thesection}{0.6em}{}
\titleformat{\subsection}{\large\bfseries\color{litblue}}{\thesubsection}{0.6em}{}
\usepackage{tcolorbox}
\usepackage{fancyhdr}
\pagestyle{fancy}
\fancyhf{}
\fancyhead[L]{\small\color{litblue}Froth — Apuntes de ML}
\fancyhead[R]{\small\thepage}
\renewcommand{\headrulewidth}{0.4pt}
\usepackage[colorlinks=true,linkcolor=litblue,urlcolor=litblue]{hyperref}

% Cabecera de cada apunte: \cabecera{numero}{titulo}
\newcommand{\cabecera}[2]{%
  \begin{tcolorbox}[colback=litblue,coltext=white,colframe=litblue,arc=2mm]
    {\Large\bfseries Apunte #1 — #2}\\[3pt]
    {\small Proyecto Froth · aprender Machine Learning construyendo}
  \end{tcolorbox}\vspace{0.8em}}

% Cajas temáticas
\newtcolorbox{concepto}[1]{colback=blue!4,colframe=litblue,title={Concepto: #1},arc=1.5mm}
\newtcolorbox{practica}{colback=green!5,colframe=litgreen,title={Pruébalo tú (teclado en mano)},arc=1.5mm}
\newtcolorbox{ojo}{colback=orange!8,colframe=litorange,title={Ojo / error típico},arc=1.5mm}
\newtcolorbox{resumen}{colback=gray!8,colframe=gray!60,title={En una frase},arc=1.5mm}

\newcommand{\codigo}[1]{\texttt{#1}}
 justo después de \documentclass.
\usepackage[margin=2.4cm]{geometry}
\usepackage{xcolor}
\definecolor{litblue}{RGB}{31,79,121}
\definecolor{litgreen}{RGB}{27,94,32}
\definecolor{litorange}{RGB}{230,81,0}
\usepackage{parskip}
\usepackage{titlesec}
\titleformat{\section}{\Large\bfseries\color{litblue}}{\thesection}{0.6em}{}
\titleformat{\subsection}{\large\bfseries\color{litblue}}{\thesubsection}{0.6em}{}
\usepackage{tcolorbox}
\usepackage{fancyhdr}
\pagestyle{fancy}
\fancyhf{}
\fancyhead[L]{\small\color{litblue}Froth — Apuntes de ML}
\fancyhead[R]{\small\thepage}
\renewcommand{\headrulewidth}{0.4pt}
\usepackage[colorlinks=true,linkcolor=litblue,urlcolor=litblue]{hyperref}

% Cabecera de cada apunte: \cabecera{numero}{titulo}
\newcommand{\cabecera}[2]{%
  \begin{tcolorbox}[colback=litblue,coltext=white,colframe=litblue,arc=2mm]
    {\Large\bfseries Apunte #1 — #2}\\[3pt]
    {\small Proyecto Froth · aprender Machine Learning construyendo}
  \end{tcolorbox}\vspace{0.8em}}

% Cajas temáticas
\newtcolorbox{concepto}[1]{colback=blue!4,colframe=litblue,title={Concepto: #1},arc=1.5mm}
\newtcolorbox{practica}{colback=green!5,colframe=litgreen,title={Pruébalo tú (teclado en mano)},arc=1.5mm}
\newtcolorbox{ojo}{colback=orange!8,colframe=litorange,title={Ojo / error típico},arc=1.5mm}
\newtcolorbox{resumen}{colback=gray!8,colframe=gray!60,title={En una frase},arc=1.5mm}

\newcommand{\codigo}[1]{\texttt{#1}}

\begin{document}
\cabecera{24}{Qué es un gap (y cómo se mide de verdad)}

\section{Por qué importan}

Toda tesis debe aportar algo nuevo, y ``nuevo'' casi siempre significa \textbf{trabajar donde
otros no han trabajado}. El problema: la frase ``research gap'' se usa mucho y se mide poco —
suele ser una corazonada del supervisor. El diferenciador de Froth es convertir el gap en
una \textbf{señal medible sobre los datos}, con evidencia que puedas citar.

\section{Los 4 tipos de gap medibles}

\begin{concepto}{1. El SILO (el más valioso)}
Dos comunidades que estudian cosas \textbf{muy parecidas} pero \textbf{no se citan}: trabajan
de espaldas. Se mide restando dos matrices que ya tenemos: similitud semántica entre clusters
(alta) vs.\ citas cruzadas (bajas). \textbf{Tu corpus lo tiene}: geología $\leftrightarrow$
flotación $=$ similitud 0.69 con \textbf{una sola cita cruzada} (Apunte 21). Dos mundos
hablando del mismo granito sin hablarse. Quien conecte esos mundos, aporta.
\end{concepto}

\begin{concepto}{2. La ZONA RALA}
Regiones vacías del mapa \emph{entre} nubes densas: combinaciones de temas que nadie tocó.
Se mide con densidad espacial en el espacio de embeddings (¿cuántos papers viven entre el
centroide de A y el de B?).
\end{concepto}

\begin{concepto}{3. El PUENTE ESCASO}
Entre dos clusters solo existen 1--2 papers puente (la betweenness del Apunte 20 concentrada
en poquitos nodos). Si cruzar ese puente produce valor y casi nadie lo ha cruzado, hay sitio
para más viajeros.
\end{concepto}

\begin{concepto}{4. El TEMA DORMIDO}
Un subtema cuyos papers son todos viejos: conocimiento que nadie revisita con métodos
modernos. Se mide con la distribución de años por cluster (¿cuándo fue el último paper?).
\end{concepto}

\section{La advertencia: mina o desierto}

\begin{ojo}
Un hueco puede estar vacío por dos razones: porque \textbf{nadie lo exploró} (una mina
esperando) o porque \textbf{no hay nada que encontrar} (un desierto — a veces la combinación
de temas es estéril por buenas razones físicas o económicas). \textbf{Ningún algoritmo
distingue una de otra.} Por eso Froth \emph{propone candidatos con evidencia} y el experto
decide. Desconfía de cualquier herramienta que prometa ``gaps automáticos garantizados'';
la nuestra promete candidatos honestos con sus números.
\end{ojo}

\section{Tu tesis ya es la prueba del concepto}

El título de la tesis (flotación de micas de Li \textbf{y su impacto en} la recuperación de
by-products del granito de Beauvoir) explota \emph{exactamente} el silo \#1: une el cluster
de flotación con el de geología. Fareed encontró ese gap a ojo y con experiencia; Froth lo
encontró con números — similitud alta, una cita cruzada, y el paper de caracterización de
Beauvoir como puente solitario (Apunte 20). Esa doble confirmación (experto + datos) es la
validación más fuerte que puede tener el método.

\section{Lo que viene (5.1)}

Implementar las cuatro señales en \codigo{gaps.py} $\rightarrow$ \codigo{find\_gaps()}
devolverá una tabla rankeada: cada gap candidato con su tipo, su evidencia (similitud, citas,
densidad, años) y los papers frontera que ya lo rondan. Después (5.2), el borrador de review
citará esos gaps con números — no con corazonadas.

\begin{resumen}
Un gap medible tiene 4 formas: silos (parecidos que no se citan), zonas ralas, puentes
escasos y temas dormidos. La herramienta propone candidatos con evidencia; el experto
distingue minas de desiertos. Tu propia tesis es el silo geología-flotación de Beauvoir,
ahora confirmado con números.
\end{resumen}

\end{document}
